\section{The Electron Emerges}

The electron enters only after the grammar has exhausted cheaper readings.
An unanchored drift cannot be charged, because the boundary has not yet
forbidden it. A first residual cannot be charged at stationarity, because the
stationary reader has forced it to vanish. A merely positive activity cannot
yet be an electron, because magnitude alone has no orientation. The electron
appears when the second variation is read by an action that discriminates the
coupled residue rather than letting it telescope into silence.

The discriminating action is the action that refuses the wrong cancellation.
A telescoping action reads adjacent debts so that the left contribution is
answered by the right and the boundary closes without a signed remainder.
Such an action is necessary in the earlier chapters because it explains how
stationarity can be certified. But the pair problem is different. A pair kick
couples two readings that cannot both be absorbed as independent left and
right debts. The grammar must ask whether the coupled difference is merely
the sum of its separated parts or whether a residue remains after those
parts have been accounted for.

The electron is the remaining reading. The grammar compares the coupled
cubic difference with the left and right readings taken separately. If the
coupled reading were no more than those separated readings, the pair would
not carry a new charge. But the mixed residue is not permitted to vanish by
renaming. It survives the subtraction of the separated parts. Because the
first variation has already been driven to zero, this surviving residue is
read at second order. Because the baseline is flat, the orientation of that
second-order residue is fixed. The grammar reads the charge as \(-1\).

This sign is not an ornament attached to a particle after the fact. It is the
orientation of the discriminated residue. The flat path supplies the native
baseline: the comparison in which no tilt has been introduced, no loop has
changed the orientation, and no external convention has reversed the sign.
Relative to that baseline, the coupled residue points against the positive
orientation. The grammar therefore reads electron, not because a name has
been announced, but because the only admissible charged reading of the flat
pair is \(-1\).

The word charge should be handled with the same discipline as the word
energy in the previous chapter. Energy became readable only after the
anchored operator could detect zero activity. Charge becomes readable only
after the baseline can discriminate orientation. A magnitude without a
baseline has size but no sign. A sign without the second variation has
orientation but no activity to carry it. The electron combines the two: a
second-variation activity whose discriminating action reads negative
orientation over the flat path.

This also explains why the electron emerges from the invariant rather than
beside it. The grammar has not introduced a second measure. It has taken the
single positive invariant and asked how it reads when the pair residue cannot
be decomposed into separated debts. The invariant supplies the activity; the
baseline supplies orientation; the discriminating action supplies the refusal
of false cancellation. The result is a charged reading of the same invariant.

The flatness of the path matters because it forbids the positron reading at
this stage. Over the native baseline, a \(+1\) charge would require either a
tilt in the comparison or a reversal of the sign convention. Neither is
available inside the flat reading. The grammar therefore does not merely
permit the electron reading. It forbids the opposite reading until the
baseline changes. The electron is forced because every other charged
interpretation would spend a distinction the flat grammar has not granted.

A finite coincidence example can illuminate the sign without supplying the
argument. When a detector later reads two back-to-back photons inside a
finite energy window, it treats a prior electron-positron encounter as a
cancellation of signed charges. The cancellation is visible only because the
opposed signs have already been distinguished. In the present section the
grammar is earlier than that detector story. It is establishing the native
negative sign that such a later cancellation can use.

Thus the electron is not an assumed inhabitant of the ledger. It is the
grammar's first charged reading of the single invariant. The second
variation carries activity. The pair kick exposes a mixed residue. The
discriminating action refuses to quotient that residue away. The flat
baseline fixes the orientation. Those obligations close on the charge
\(-1\), and the name electron is the reading of that closure.

